%%%%%%%%%%%%%%%%%%%%%%%%%%%%%%%%%%%%%%%%%%%%%%%%%%%%%%%%%%%%%%%%%%%%%%
% How to use writeLaTeX: 
%
% You edit the source code here on the left, and the preview on the
% right shows you the result within a few seconds.
%
% Bookmark this page and share the URL with your co-authors. They can
% edit at the same time!
%
% You can upload figures, bibliographies, custom classes and
% styles using the files menu.
%
%%%%%%%%%%%%%%%%%%%%%%%%%%%%%%%%%%%%%%%%%%%%%%%%%%%%%%%%%%%%%%%%%%%%%%

\documentclass[12pt]{article}

\usepackage{sbc-template}
\usepackage{listings}
\usepackage{graphicx,url}
\usepackage{amsmath}
%\usepackage[brazil]{babel}   
\usepackage[utf8]{inputenc}  

     
\sloppy

\title{Análise da Lógica Ternária e Eficiência Computacional}

\author{João Ricardo Pereira Marques\inst{1} }


\address{Bacharelado em Ciência da Computação - Universidade Tecnológica Federal do Paraná\\
  Campus Campo Mourão - Campo Mourão - Paraná - Brasil
  \email{joaoricardomarques@alunos.utfpr.edu.br}
}

\begin{document} 
\maketitle
\begin{abstract}
	teste
\end{abstract}
\begin{resumo}
	teste
\end{resumo}

\section{Introdução}
A computação ternária é uma alternativa direta ao modelo binário, limitado a apenas 0 e 1, enquanto o sistema ternário conta com três estados ("0,1,2"ou "-1,0,1") assim sento matematicamente mais eficiente na codificação de informação em um único \textbf{trit}
\section{Fundamentação}
\subsection{Conversão de ternário para decimal}

\begin{equation}
	N_{10} = (d_n \cdot 3^n)+(d_{n-1} \cdot 3^{n-1})+....+(d_1 \cdot 3^1)+(d_0 \cdot 3^0)
\end{equation}

\subsection{Conversão de decimal para ternário}
\begin{equation}
	\begin{split}
		N   &= 3 \cdot q_0 + r_0 \\
		q_0 &= 3 \cdot q_1 + r_1 \\
		q_1 &= 3 \cdot q_2 + r_2 \\
		&\vdots \\
		q_{n-1} &= 3 \cdot q_n + r_n
	\end{split}
\end{equation}
\section{Operações lógicas e aritméticas}
\subsection{NOT}
\centering
\begin{tabular}{|c|c|c|c|}
	\hline
	\textbf{A} & PNOT \textbf{A}& SNOT \textbf{A}& NNOT \textbf{A}\\\hline
	0&2&2&2\\\hline
	1&2&1&0\\\hline
	2&0&0&0\\\hline
\end{tabular}
\subsection{AND}
\centering
\begin{tabular}{|c|c|c|}
	\hline
	\textbf{A} & \textbf{B} & \textbf{A} AND \textbf{B}\\\hline
	0&0&0\\\hline
	0&1&0\\\hline
	0&2&0\\\hline
	1&0&0\\\hline
	1&1&1\\\hline
	1&2&1\\\hline
	2&0&0\\\hline
	2&1&1\\\hline
	2&2&2\\\hline
\end{tabular}
\subsection{OR}
\centering
\begin{tabular}{|c|c|c|}
	\hline
	\textbf{A} & \textbf{B} & \textbf{A} OR \textbf{B}\\\hline
	0&0&0\\\hline
	0&1&1\\\hline
	0&2&2\\\hline
	1&0&1\\\hline
	1&1&1\\\hline
	1&2&2\\\hline
	2&0&2\\\hline
	2&1&2\\\hline
	2&2&2\\\hline
\end{tabular}
\subsection{XOR}
\centering
\begin{tabular}{|c|c|c|}
	\hline
	\textbf{A} & \textbf{B} & \textbf{A} XOR \textbf{B}\\\hline
	0&0&0\\\hline
	0&1&1\\\hline
	0&2&2\\\hline
	1&0&1\\\hline
	1&1&0\\\hline
	1&2&1\\\hline
	2&0&2\\\hline
	2&1&1\\\hline
	2&2&0\\\hline
\end{tabular}
\subsection{XNOR}
\centering
\begin{tabular}{|c|c|c|}
	\hline
	\textbf{A} & \textbf{B} & \textbf{A} XNOR \textbf{B}\\\hline
	0&0&2\\\hline
	0&1&1\\\hline
	0&2&0\\\hline
	1&0&1\\\hline
	1&1&2\\\hline
	1&2&1\\\hline
	2&0&0\\\hline
	2&1&1\\\hline
	2&2&2\\\hline
\end{tabular}
\subsection{Soma}
\centering
\begin{tabular}{|c|c|c|c|c|}
	\hline
	recebe-um&\textbf{A} & \textbf{B} & \textbf{A} + \textbf{B} & vai-um\\\hline
	0&0&0&0&0\\\hline
	0&0&1&1&0\\\hline
	0&0&2&2&0\\\hline
	0&1&0&1&0\\\hline
	0&1&1&2&0\\\hline
	0&1&2&0&1\\\hline
	0&2&0&2&0\\\hline
	0&2&1&0&1\\\hline
	0&2&2&1&1\\\hline
	1&0&0&1&0\\\hline
	1&0&1&2&0\\\hline
	1&0&2&0&1\\\hline
	1&1&0&2&0\\\hline
	1&1&1&0&1\\\hline
	1&1&2&1&1\\\hline
	1&2&0&0&1\\\hline
	1&2&1&1&1\\\hline
	1&2&2&2&1\\\hline
\end{tabular}
\subsection{Subtração}
\centering
\begin{tabular}{|c|c|c|c|c|}
	\hline
	empresta-um & \textbf{A} & \textbf{B} & \textbf{A} - \textbf{B}& empresta-um \\\hline
	0 & 0 & 0 & 0 & 0 \\\hline
	0 & 0 & 1 & 2 & 1 \\\hline
	0 & 0 & 2 & 1 & 1 \\\hline
	0 & 1 & 0 & 1 & 0 \\\hline
	0 & 1 & 1 & 0 & 0 \\\hline
	0 & 1 & 2 & 2 & 1 \\\hline
	0 & 2 & 0 & 2 & 0 \\\hline
	0 & 2 & 1 & 1 & 0 \\\hline
	0 & 2 & 2 & 0 & 0 \\\hline
	1 & 0 & 0 & 2 & 1 \\\hline
	1 & 0 & 1 & 1 & 1 \\\hline
	1 & 0 & 2 & 0 & 1 \\\hline
	1 & 1 & 0 & 0 & 0 \\\hline
	1 & 1 & 1 & 2 & 1 \\\hline
	1 & 1 & 2 & 1 & 1 \\\hline
	1 & 2 & 0 & 1 & 0 \\\hline
	1 & 2 & 1 & 0 & 0 \\\hline
	1 & 2 & 2 & 2 & 1 \\\hline
\end{tabular}
\subsection{Multiplicação}

\subsection{Divisão}

\section{Implementação de Portas}
\subsection{Gatilho de Schmitt}

\subsection{Porta SNOT}

\subsection{Porta PNOT}

\subsection{Porta NNOT}

\subsection{Porta AND}

\subsection{Porta OR}

\subsection{Porta SNAND}

\subsection{Porta PNAND}

\subsection{Porta NNAND}

\subsection{Porta SNOR}

\subsection{Porta PNOR}

\subsection{Porta NNOR}

\section{Implementação de Circuito }

\subsection{Somador}

\subsection{Subtrator}

\subsection{Multiplicador}

\subsection{Divisor}

\subsection{Comparador}
\bibliographystyle{sbc}
\bibliography{sbc-template}

\end{document}
